\documentclass[10pt,a4paper]{amsart}
\usepackage[margin=19mm]{geometry}
\usepackage{amsmath,amssymb,mathtools}
\usepackage[scr=rsfs,cal=euler]{mathalfa}
\usepackage{xfrac}
\usepackage[unicode,hidelinks,bookmarks=false]{hyperref}
\newcommand{\ie}{i.e.\ }
\newcommand{\cf}{cf.\ }
\newcommand{\resp}{respectively\ }
\newcommand{\Subsection}{\subsection}
\providecommand{\nobreakdash}{\nobreak\hbox{-}}
\setlength{\emergencystretch}{2em}
\multlinegap0pt
\usepackage{bm}
\usepackage{bbm}
\usepackage{dsfont}
\usepackage{xspace}
\usepackage{cancel}
\usepackage{mathtools}
\usepackage{amsthm}
\makeatletter
\@ifundefined{thm}{\newtheorem{thm}{Thm}}{}
\@ifundefined{lem}{\newtheorem{lem}[thm]{Lemma}}{}
\makeatother
\DeclareMathOperator{\Id}{Id}
\DeclareMathOperator{\pt}{pt}
\def\bC{\mathbb{C}}
\def\bH{\mathbb{H}}
\def\bZ{\mathbb{Z}}
\title[Modularity of $d$-elliptic loci with level structure]{Modularity of $d$-elliptic loci with level structure}
\author{Fran\c{c}ois Greer, Carl Lian, Naomi Sweeting}
\date{Source: arXiv:2411.00957v4 (CC BY 4.0)}
\begin{document}
\maketitle

\noindent\emph{Source and licence.} Modularity of $d$-elliptic loci with level structure. Version arXiv:2411.00957v4 (CC BY 4.0), distributed under the Creative Commons Attribution 4.0 International licence, as recorded in the arXiv licence metadata for this exact version and shown on the abstract page https://arxiv.org/abs/2411.00957v4. Full rights notice: licenses/arxiv-2411.00957-CC-BY-4.0.txt.\par\smallskip

\noindent\textit{Editorial scope.} Counted unit: the Lem whose source label is \texttt{lem:B\_vanishing} in the article (source0.tex lines 559--565, source label \texttt{lem:B\_vanishing}), delivered together with the article's own complete original proof of it (source0.tex lines 566--618, one \texttt{proof} environment of 53 lines). No mathematics is written, repaired or re-derived: statement and proof are the article's own text line for line. Required context retained uncounted: none. Every definition, display and label of those lines is retained unchanged. Omitted from this package and not redistributed: 1--558, 619--1242. Nothing inside a retained excerpt is omitted or altered: every retained body is delivered exactly as the article's lines are written, so no omission receipt of any kind accompanies this package. Citations: the retained excerpts carry no citation command at all, so no bibliography entry of the article is transcribed and no reference is redistributed. Reference adaptation: none was needed and none was made; every reference inside the delivered unit resolves inside it, so no unresolved reference or citation is produced. Printed slips kept literal: no printed word, symbol, hypothesis or number of the counted statement or of its proof was altered. Layout and class: the article's own class is \texttt{amsart} with packages including nag, fullpage, url, amssymb, enumerate, colonequals, xy, comment, graphicx, fontenc, color, microtype; the wrapper is the lane's pinned \texttt{amsart} editorial wrapper at 10pt on A4, which restates the theorem-like environments the retained text needs and adds the article's own macro definitions that the retained text needs (\texttt{\textbackslash Id}, \texttt{\textbackslash bC}, \texttt{\textbackslash bH}, \texttt{\textbackslash bZ}, \texttt{\textbackslash pt}), with extra wrapper packages (bm, bbm, dsfont, xspace, cancel, mathtools). The delivered numbering is the unit's own. Assets: no retained span contains a figure environment, a graphics-inclusion command, a PDF-page inclusion, a table or a TikZ picture; no image file is copied into this package, the package declares no asset dependency and no image file is redistributed with it. Licence: version arXiv:2411.00957v4 of this article is distributed under the Creative Commons Attribution 4.0 International licence, as recorded in the arXiv licence metadata for this exact version and shown on the abstract page https://arxiv.org/abs/2411.00957v4. A full rights notice is delivered at licenses/arxiv-2411.00957-CC-BY-4.0.txt. Characters: every retained excerpt is pure ASCII, so the delivered engine, XeLaTeX with its default Latin Modern text font, reports no missing character.
\begin{lem}\label{lem:B_vanishing}
Given a pair $(\tau,\tau')\in \bH\times \bH_g$, there exists a degree $d$ map $E_\tau \to A_{\tau'}$ between the associated abelian varieties if and only if 
\begin{equation*}
    (B \tilde{\tau})^T\cdot \tilde{\tau}'=0\in\bC^g
\end{equation*}
for some matrix $B\in M_{2g\times 2,d}$. Here, $\tilde{\tau} = \begin{pmatrix} \tau \\ \Id \end{pmatrix}$ denotes the Siegel augmentation.
\end{lem}

\begin{proof}
Let $h: E \to A$ be a degree $d$ map. Let $B_h\in M_{2g\times 2,d}$ be the matrix for 
$$h_*:H_1(E) \to H_1(A)$$
with respect to symplectic bases $\langle\alpha,\beta\rangle$ for $H_1(E)$ and $\langle\alpha_j,\beta_j\rangle_{j=1,\ldots,g}$ for $H_1(A)$. The graph $\Gamma_h \subset E \times A$ has homology class
\begin{equation*}
\alpha\times h_*\beta - \beta\times h_*\alpha + E\times [\pt] + [\pt]\times h_*[E]\in H_2(E\times A,\bZ).
\end{equation*}

If $\omega\in H^{1,0}(E)$ and $\omega_j\in H^{1,0}(A)$ are the normalized holomorphic 1-forms, then their external wedge product $\omega \wedge \omega_j\in H^{2,0}(E\times A)$ integrates to 0 on any algebraic curve, so on $\Gamma_h$ in particular. In terms of the decomposition above, this implies that
\begin{equation*}
\int_\alpha \omega \int_\beta h^*\omega_j - \int_\beta \omega \int_\alpha h^*\omega_j=0
\end{equation*}
for $j=1,2,\dots,g$. 

The Siegel augmented period matrices are given by
\begin{align*}
    \tilde{\tau} &= \begin{pmatrix} \int_\alpha \omega \\ \int_\beta \omega  \end{pmatrix},\\
    \tilde{\tau}' &= \begin{pmatrix} \int_{\alpha_i} \omega_j \\ \int_{\beta_i} \omega_j  \end{pmatrix}.
\end{align*}
Multiplying $\tilde{\tau}'$ by $B_h^T$ on the left has the effect of replacing $\alpha_i$ and $\beta_j$ with the pushforwards of $\alpha$ and $\beta$ under $h$:
\begin{equation*}
B_h^T \tilde{\tau}' = 
\begin{pmatrix} 
\int_{h_*\alpha} \omega_j \\ 
\int_{h_*\beta} \omega_j  \end{pmatrix} 
=
\begin{pmatrix} 
\int_{\alpha} h^*\omega_j \\ 
\int_{\beta} h^*\omega_j 
\end{pmatrix}
.
\end{equation*}
Now, we also have
\begin{equation*}
J_2^{-1}\tilde{\tau} 
=
\begin{pmatrix}
- \int_\beta \omega \\ 
\int_\alpha \omega 
\end{pmatrix}
\end{equation*}
and the vanishing above is equivalent to 
\begin{equation*}
    \tilde{\tau}^T J_2 B_h^T \tilde{\tau}'= (B_hJ_2^{-1}\tilde{\tau})^T\cdot \tilde{\tau'} =0.
\end{equation*}
Taking $B=B_hJ_2^{-1}$ yields the first direction of the Lemma.

Conversely, given $B\in M_{2g\times 2,d}$, define $B_h=BJ_2$, and a linear subtorus 
\begin{equation*}
    \Gamma\subset E\times A=\bC/\Lambda_\tau\times \bC^g/\Lambda'_{\tau'}
\end{equation*}
by $\Gamma=\{(z,B_hz)\mid z\in\bC\}$, where we extend the symplectic bases of the lattices $\Lambda_\tau,\Lambda'_{\tau'}$ to $\bC,\bC^g$, respectively. Reversing the previous calculation, the vanishing $(B \tilde{\tau})^T\cdot \tilde{\tau}'=0$ implies that integrals of holomorphic 2-forms on $\Gamma$ vanish, which in turn implies that $\Gamma\cong E$ is a complex subtorus. Post-composing with the projection to $A$ gives the desired map $h:E\to A$.
\end{proof}

\end{document}
